\chapter{Where Clauses, Identity and Logs}\label{ch:where}

A session type is abstract: it says \emph{GetImage, then Image or NotFound}. A
running system sends bytes that carry ids, amounts and tokens. Where clauses are the
bridge. They state, in the protocol, the relations between field values that make a
sequence of messages one conversation and not a coincidence of labels. This chapter
shows the three jobs they do (data conditions, identity and correlation) and then the
division of labour the rest of the book relies on: CTL decides the order of messages,
and where clauses decide the values that implement that order.

\section{Conditions on data}

A where clause on a message is a condition over that message's fields and the fields
of earlier messages it names. PlaceOrder requires \texttt{p.amount = q.price}. Given a
log, \texttt{session\_conform} replays it against the protocol and evaluates every
where clause on the values the log carries. Here the client pays 12 against a quote
of 10:

\begin{lstlisting}[language={},basicstyle=\ttfamily\footnotesize,breaklines=true]
{"from":"Client","id":"e1","label":"CreateOrder","session":"7","to":"Shop"}
{"cause":"e1","fields":{"price":10},"from":"Shop","id":"e2","label":"Quote","session":"7","to":"Client"}
{"cause":"e2","fields":{"amount":12},"from":"Client","id":"e3","label":"Pay","session":"7","to":"Shop"}
\end{lstlisting}

\begin{lstlisting}[language={},basicstyle=\ttfamily\footnotesize,breaklines=true]
PlaceOrder: 1 of 1 sessions break the protocol; first: session 7, event e3: Pay from Client to Shop breaks its condition: where p.amount = q.price is false: p.amount = 12, q.price = 10
\end{lstlisting}


The order of the three messages is one the protocol allows. The data is not. A
where clause is therefore a second, independent check on a run: the first asks whether
the labels form a path of the protocol, the second whether the values along that path
satisfy its conditions.

\begin{definition}[Conformance]\label{def:conform}
A log $e_1 e_2 \ldots e_k$ \emph{conforms} to a protocol when (control) the sequence of
its labels, with their causes, is a path from the start of the protocol's state graph,
and (data) for each $e_i$ every where clause on its message holds of the field values
of $e_i$ and of the earlier events its clause names, as the causes and bound ids
identify them.
\end{definition}

\section{Identity: telling sessions apart}

A real log interleaves many sessions, and most logs do not say which session an event
belongs to. LeanFM derives sessions from the protocol: a \texttt{cause} joins an answer
to its request, and fields the protocol marks as identities join the events that carry
the same value. Two image requests interleaved, with no session ids at all:

\begin{lstlisting}[language={},basicstyle=\ttfamily\footnotesize,breaklines=true]
{"fields":{"id":"cat"},"from":"Browser","id":"a1","label":"GetImage","to":"Server"}
{"fields":{"id":"dog"},"from":"Browser","id":"b1","label":"GetImage","to":"Server"}
{"cause":"b1","from":"Server","id":"b2","label":"NotFound","to":"Browser"}
{"cause":"a1","fields":{"bytes":"...","id":"cat"},"from":"Server","id":"a2","label":"Image","to":"Browser"}
\end{lstlisting}

\begin{lstlisting}[language={},basicstyle=\ttfamily\footnotesize,breaklines=true]
Images: 4 events in 2 sessions follow the protocol

derived sessions:
  chain from a1: events [0,3]
  chain from b1: events [1,2]
\end{lstlisting}


The causes are enough. In general the derivation is a union-find over the log: each
cause joins two events, each declared key joins the events that carry its value, and
each resulting group is checked as its own run. Two events are never joined because
their payloads happen to be equal; only the protocol's marks join them.

\paragraph{Marks.} A field's type says what it does for identity. A field of type
\texttt{Fresh String} mints a new session id; \texttt{Link String} names an existing
one; \texttt{Header}, \texttt{PathParam} and \texttt{QueryParam} say where it travels
on the wire. A protocol that is \texttt{given sessionId} runs inside a session minted
elsewhere; including it \texttt{with sessionId} binds that field to the caller's.

\section{Sessions that fork and join}

GoServer is a multiplayer Go server in this repository. A person signs in, which mints
a session; a game runs as a session of its own; a seat at a game belongs to both the
person and the game.

\begin{lstlisting}[language={},basicstyle=\ttfamily\footnotesize,breaklines=true]
protocol GameP (Person, Server, Bot) given sessionId {
  g: CreateGame from Person to Server route "POST /api/create_game";
  choice {
    c: GameCreated from Server to Person cause g status 200 where c.requestId = g.requestId;
    include Challenge Person Server with sessionId, gameId;
    rec Turns {
      choice {
        include Turn Server Bot;
        continue Turns;
      } or {
        Finished at Server;
      }
    }
  } or {
    x: RequestRejected from Server to Person cause g status 200 where x.requestId = g.requestId;
  }
}
\end{lstlisting}

\begin{lstlisting}[language={},basicstyle=\ttfamily\footnotesize,breaklines=true]
-- from projects/GoServer/source/cmd/demo/static/session/goserver.leanfm
protocol SeatP (Person, Server) given sessionId gameId {
  include TakeSeatP Person Server with sessionId, gameId;
  parallel {
    sessions { include PlayMoveP Person Server with sessionId, gameId; }
  } and {
    sessions { include SetAutopilotP Person Server with sessionId, gameId; }
  } and {
    sessions {
      choice {
        include ProposeMergeP Person Server with sessionId, gameId;
      } or {
        include RespondMergeP Person Server with sessionId, gameId;
      }
    }
  } and {
    sessions { include GameStateP Person Server with sessionId, gameId; }
  }
}
\end{lstlisting}


The game is \texttt{given sessionId}, the person's, so it is a child of that person's
session. A seat is \texttt{given sessionId gameId}: it names two sessions that are not
parent and child, so it is a \emph{join}. LeanFM reads that shape from the text alone
and checks a real log of two people playing against it:

\begin{lstlisting}[language={},basicstyle=\ttfamily\footnotesize,breaklines=true]
GoServer: 24 events in 7 sessions follow the protocol; 3 still running

chain from i1: 6 events
chain from i2: 6 events
chain from c1: 4 events, child of chain from i1
chain from t1: 2 events, joins chain from i1, chain from c1
chain from t2: 2 events, joins chain from i2, chain from c1
chain from q1: 2 events, joins chain from i1, chain from c1
chain from q2: 2 events, joins chain from i2, chain from c1
\end{lstlisting}


The log carries no session ids. LeanFM found the two people's sessions, the game as a
child of the first person's session, and each seat and move as a join of a person and
the game. Three sessions are still running when the log ends, which is why the call
passes \texttt{allowIncomplete}.

\section{Where clauses implement prev and cause}

In a diagram, a message's place in the conversation is drawn: an arrow below another
arrow. On the wire nothing is drawn. Each message has to carry enough to say where it
belongs, and that is what where clauses state:

\begin{itemize}
\item \texttt{cause g} with \texttt{where i.id = g.id} says the image answers
  \emph{this} request. On the wire the answer repeats the request's identity
  (a request id, a stream id, the tuple of a TCP connection).
\item A message's \emph{prev} is the message before it in the same thread. A message
  that carries a pointer to its predecessor carries its own place in the protocol: its
  typestate is its predecessor's typestate after one step.
\end{itemize}

No where clause needs to mention prev explicitly. The message itself carries the
pointer, and the log reader follows it. Chapter~\ref{ch:constructs} shows why this
matters for checking: with prev, a typestate never has to be inferred from what may
happen later; it is read off the chain. That one property turns the cost of checking
from quadratic to linear.

\section{The division of labour}

Part~\ref{part:ctl} reduces the protocol's \emph{control}, the order and branching of
labels, to CTL. Values are abstracted: the model knows that an Image follows a
GetImage, not which image. Where clauses carry the values. The two meet in
Definition~\ref{def:conform}: a run conforms when its labels form a path of the model
the CTL rules characterise, and its values satisfy the where clauses along that path.

A where clause that selects a branch, such as \texttt{where o.total >= 100} on one
branch and \texttt{o.total < 100} on the other, is drawn as that branch's guard. In the
CTL of the control it appears as an ordinary choice: both branches are possible as far
as order goes, and the guard decides which one a given run takes. The guard is checked
on the values; the order is checked by CTL.
